% Propietario conservado: /Users/ruben/Documents/New project/output/EXTREMA_MEDIA_RAZON_RECIPROCIDAD_ES_EN_20260918/ARTICULO/ES/source/nuclear/02.tex
\chapter{Conservation of observables and the complete record}
\label{x_reciprocidad:nuclear:2}
\section{Balance of a conserved observable}

Let \(\mathcal H,\mathcal K\) be realizations of a readout space and its enriched extension. Let \(J:\mathcal H\to\mathcal K\) be an isometry, \(P=JJ^*\), \(U:\mathcal K\to\mathcal K\) a bounded transport, and \(\widetilde A\) a bounded self-adjoint observable. Suppose

\[
U^*\widetilde A U=\widetilde A,\qquad [\widetilde A,P]=0.
\]

We define

\[
A=J^*\widetilde A J,\qquad T=J^*UJ,
\qquad \eta=(I-P)UJ.
\]

\begin{proposition}
\label{x_reciprocidad:nuc02:balance}
The charge decomposes exactly between readout and complement:

\begin{equation}
\boxed{A=T^*AT+\eta^*\widetilde A\eta.}\label{x_reciprocidad:nuc02:eq:1}
\end{equation}
\end{proposition}

\begin{proof}
The identity \(UJ=JT+\eta\) is an orthogonal decomposition. Expanding \(J^*U^*\widetilde A UJ\), the cross terms vanish because \(\widetilde A\) reduces the subspaces \(P\mathcal K\) and \((I-P)\mathcal K\). The left-hand term is \(A\); the two diagonal terms are those in \eqref{x_reciprocidad:nuc02:eq:1}. This also proves the identity of sesquilinear forms for two distinct states. If \(\widetilde A\geq0\), both summands are positive. For \(\widetilde A=I\), the conservation hypothesis is \(U^*U=I\).
\end{proof}

If the readout mixes sectors of the observable, the balance additionally retains

\[
T^*J^*\widetilde A\eta+\eta^*\widetilde A JT.
\]

These terms represent interference between the two spaces carrying the charge. The reducing condition stated above ensures their vanishing for the entire decomposition. In particular realizations, they may also vanish for other reasons, for example when \(\eta=0\).

\section{Specialization to the nonadic connection}

Let \(C\) be the unitary realization of memory transport on \(\mathcal H\). On \(\mathcal H^9\), transport of the nine phases and their uniform inclusion are

\[
S_C(v_0,\ldots,v_8)=(Cv_8,v_0,\ldots,v_7),
\qquad Jv=\tfrac13(v,\ldots,v).
\]

The factor \(1/3\) follows from normalization of nine copies. The compression and its complement are

\[
T=\frac{8I+C}{9},\qquad
\eta v=\frac1{27}\bigl(8(C-I)v,-(C-I)v,\ldots,-(C-I)v\bigr).
\]

If \(A=A^*\) is bounded and \([A,C]=0\), the observable \(\widetilde A=I_9\otimes A\) satisfies the hypotheses of Proposition~\ref{x_reciprocidad:nuc02:balance}. Therefore,

\begin{equation}
\boxed{A=T^*AT+\frac8{81}(I-C)^*A(I-C).}\label{x_reciprocidad:nuc02:eq:2}
\end{equation}

The coefficients arise from the distribution among the nine components: \(8^2+8=72\), divided by \(27^2\), gives \(8/81\). The identity is the sesquilinear balance already documented in the spectral article, expressed here for a general observable.

Telescopic iteration gives, for every \(N\geq1\),

\begin{equation}
A=T^{*N}AT^N+
\sum_{j=0}^{N-1}T^{*j}\eta^*\widetilde A\eta T^j.\label{x_reciprocidad:nuc02:eq:3}
\end{equation}

For \(A=I\), the map

\[
V_Nv=(T^Nv,\eta v,\eta Tv,\ldots,\eta T^{N-1}v)
\]

is isometric. The step \(N\to N+1\) decomposes only the terminal component \(T^Nv\) into its new readout and complement, retaining every preceding record. Conservation is thus verified at any depth.

Distinction of operations. One has \(S_C^9=I_9\otimes C\). Iteration \(T^N\) describes successive compressions with storage of their complements. Enriched transport without intermediate compressions is \(S_C^N\). Both operations have defined domains and different functions; the proof preserves this distinction.

\section{Terminal sector and the unlimited-depth record}

Let \(P_{\rm fix}\) be the orthogonal projection onto \(\ker(I-C)\).

\begin{theorem}
For every unitary transport \(C\), the nonadic compression satisfies

\begin{equation}
\underset{N\to\infty}{\operatorname{s-lim}}T^N=P_{\rm fix},\qquad
\boxed{I=P_{\rm fix}+\sum_{j=0}^{\infty}T^{*j}\eta^*\eta T^j.}\label{x_reciprocidad:nuc02:eq:4}
\end{equation}

The series converges in the strong topology. For each bounded self-adjoint observable \(A\) commuting with \(C\), one also obtains

\begin{equation}
\boxed{A=P_{\rm fix}AP_{\rm fix}
+\sum_{j=0}^{\infty}T^{*j}\eta^*(I_9\otimes A)\eta T^j.}\label{x_reciprocidad:nuc02:eq:5}
\end{equation}
\end{theorem}

\begin{proof}
In the spectral representation of the unitary operator \(C\), compression corresponds to \(t(z)=(8+z)/9\). On the unit circle,

\[
|t(z)|^2=1-\frac8{81}|1-z|^2.
\]

Hence \(t(z)^N\to0\) for \(z\ne1\), while \(t(1)^N=1\). Dominated convergence with respect to each vector spectral measure proves \(T^N\to P_{\rm fix}\) strongly. The same argument holds for \(T^{*N}\). Since \(A\) is bounded, \(T^{*N}AT^N\to P_{\rm fix}AP_{\rm fix}\) strongly. Taking the limit in \eqref{x_reciprocidad:nuc02:eq:3} gives \eqref{x_reciprocidad:nuc02:eq:5}, and specializing to \(A=I\) gives \eqref{x_reciprocidad:nuc02:eq:4}.
\end{proof}

In particular,

\[
V_\infty v=(P_{\rm fix}v,\eta v,\eta Tv,\eta T^2v,\ldots)
\]

is an isometry. It preserves the complete inner product as well as the norm.

\begin{corollary}[Bilateral memory]
For the shift \(Ce_n=e_{n+1}\) on \(\ell^2(\mathbb Z)\), \(P_{\rm fix}=0\): a fixed sequence would be constant, and its only square-summable realization is zero. Consequently, the entire initial norm is distributed through the complement record. This specifies the unlimited behaviour of the bilateral memory realization present in the corpus.
\end{corollary}

\begin{corollary}[Finite calendar]
On a cycle of \(d\) positions, the uniform sequence determines a one-dimensional fixed subspace. This sector remains as the terminal component, while the complement accumulates the remainder. The finite calendar and the bilateral record therefore have different terminal components.
\end{corollary}

In the bilateral case, strong convergence may coexist with \(\|T^N\|=1\) for every \(N\). The proof assumes no uniform contraction rate. It also distinguishes this compression from the mode \((5/9)C\), whose geometric contraction is uniform and is studied separately in the sources.

\section{Preservation of symmetry and chart transport}

Let \(R\) be unitary and define \(C'=RCR^*\), \(A'=RAR^*\). The explicit compression formula gives

\begin{equation}
T(C')R=RT(C),\qquad
\eta(C')R=R^{\oplus9}\eta(C).\label{x_reciprocidad:nuc02:eq:6}
\end{equation}

\begin{proof}
The first equality follows by expanding \((8I+RCR^*)R/9\). The second follows from \((C'-I)R=R(C-I)\) in each of the nine components of \(\eta\).
\end{proof}

Iteration of \eqref{x_reciprocidad:nuc02:eq:6} intertwines the finite registers. Moreover, \(P_{\rm fix}(C')=RP_{\rm fix}(C)R^*\), by transport of the kernel of \(I-C\), so the unlimited register is also intertwined. Thus, the balance of observables is covariant when the chart, operator and observable are transported jointly. Symmetries commuting with \(C\) additionally act in the same fixed chart. This distinction allows exceptional transformations to be incorporated with their effectively transported frames, flags and orientations.

The composition in \cref{x_reciprocidad:lem:k-hadamard,x_reciprocidad:exc:entrelazador,x_reciprocidad:exc:isometria} uses the record's Hadamard matrix, \(H_{12}^2=4I\), and the incidence \(\mathcal I\) of the twelve positions in the 132 hexads:

\[
J_H=H_{12}/2,\qquad
\mathcal I^*\mathcal I=36I+30\mathbf1\mathbf1^*,\qquad
U_W=\mathcal I/6\big|_{\mathbf1^\perp}.
\]

Both normalized operators are isometric on the indicated domains. The composition \(B=U_WJ_H\) on \(W=J_H(\mathbf1^\perp)\) preserves inner products and intertwines the transported representation of \(M_{12}\) with the action on the hexads. The integer record \(K\) additionally preserves the integral inversion \(H_{12}/4\) on its congruence lattice; this inversion and isometric normalization have different functions. The proof, the APP–Witt–Leech prolongation, and their domains are developed here.

In this composition, incidence conservation has concrete operator content. Its extension through other transports uses \eqref{x_reciprocidad:nuc02:eq:6} and the effective intertwiners from the sources, respecting their respective groups and domains.

\subsection{Complete transport of the twelve coordinates of K}

Separation of the centred subspace can be completed by also retaining its mean. Let \(P_0=I-\mathbf1\mathbf1^*/12\), \(\mu(k)=\mathbf1^*k/12\), and

\[
F(k)=\left(\mu(k),\frac16\mathcal I P_0k\right)\in\mathbb R\oplus E_{\rm inc}.
\]

Equip the codomain with \(\|(\mu,y)\|^2=12|\mu|^2+\|y\|^2\). The incidence Gram operator proves

\begin{equation}
\|F(k)\|^2=\|k\|^2,\qquad
\boxed{F^{-1}(\mu,y)=\mu\mathbf1+\frac16P_0\mathcal I^*y.}\label{x_reciprocidad:nuc02:eq:6a}
\end{equation}

Indeed, \(\mathcal I^*\mathcal I P_0=36P_0\), so the inverse formula recovers \(P_0k\), while the first term recovers its mean. The action of \(M_{12}\) is trivial on the mean and permutes the hexads; hence \(F\) is equivariant. Composing \(FJ_H\) transports the normalized signed coordinates. All twelve record coordinates are thereby preserved, rather than only the centred part. The domain of this reconstruction is the dodecaphasic record; the complete history memory retains its other coordinates and its prolongation.

\subsection{Specialization to the exceptional operator of order three}

The exceptional construction produces the orthogonal operator \(g_c\) on the realization of the Leech neighbour, through the twelve APP–Witt fibres of type \(A_2\), and proves \(g_c^3=I\) and \(\ker(I-g_c)=0\). Its construction, intertwining, and lattice preservation are located in \cref{x_reciprocidad:nuclear:fibras-app-witt,x_reciprocidad:nuc04:c3}.

We now construct, as an operator composition of that output with the nine-phase chart,

\[
T_3=\frac{8I+g_c}{9},\qquad
\eta_3=(I-JJ^*)S_{g_c}J.
\]

\begin{corollary}
This composition satisfies

\begin{equation}
\boxed{T_3^*T_3=\frac{19}{27}I,\qquad
\eta_3^*\eta_3=\frac8{27}I,\qquad
\|T_3^Nv\|^2=\left(\frac{19}{27}\right)^N\|v\|^2.}\label{x_reciprocidad:nuc02:eq:6b}
\end{equation}
\end{corollary}

\begin{proof}
The vector \((I+g_c+g_c^2)v\) is fixed by \(g_c\); the absence of fixed vectors implies \(I+g_c+g_c^2=0\). Since \(g_c^*=g_c^2\), one has \(g_c+g_c^*=-I\). Substitution gives

\[
T_3^*T_3=\frac{65I+8(g_c+g_c^*)}{81}
=\frac{57}{81}I=\frac{19}{27}I.
\]

Balance \eqref{x_reciprocidad:nuc02:eq:2} with \(A=I\) gives the second coefficient. Repeated application of the norm identity gives the power for every \(N\).
\end{proof}

The unlimited record is isometric, and its terminal component tends to zero at the explicit rate in \eqref{x_reciprocidad:nuc02:eq:6b}. Here, the exceptional structure determines the two complementary coefficients; both arise from the polynomial relation of \(g_c\) and the nine phases. This is a mathematical composition developed in this research. Its definition neither identifies \(g_c\) with \(\Gamma_9\) nor asserts that \(S_{g_c}\) is the canonical temporal operator of TPK: it preserves the distinction between exceptional prolongation and temporal holonomy while composing them explicitly.

\section{History refinement and compatibility of balances}

For enriched histories of depth \(n\), let \(p_h(\varepsilon)\) be the probability of a surviving extension, normalized by \(\sum_\varepsilon p_h(\varepsilon)=1\). Canonical refinement on the history bases is

\[
R_n\delta_h=\sum_{\varepsilon\in\operatorname{Out}(h)}
\sqrt{p_h(\varepsilon)}\,\delta_{h\varepsilon}.
\]

Extensions of different histories have different prefixes and disjoint supports. Hence \(R_n^*R_n=I\). This is the isometric realization of projective consistency of the measure: the coefficients arise from the survival weights already constructed.

To transport \eqref{x_reciprocidad:nuc02:eq:1} between depths, inclusions, projections, transports, and observables are preserved jointly. With refinement isometries \(R_n,\widetilde R_n\), the required squares are

\[
J_{n+1}R_n=\widetilde R_nJ_n,\quad
P_{n+1}\widetilde R_n=\widetilde R_nP_n,\quad
U_{n+1}\widetilde R_n=\widetilde R_nU_n,\quad
\widetilde A_{n+1}\widetilde R_n=\widetilde R_n\widetilde A_n.
\]

These squares imply \(T_{n+1}R_n=R_nT_n\) and \(\eta_{n+1}R_n=\widetilde R_n\eta_n\). With uniform bounds, the identities pass to the inductive limit of the realizations. Histories, by contrast, form the projective system of prefixes; the opposite directions of the two types of maps remain distinguished. The measure, Gram, and passage-to-the-limit proofs are included.

\section{Variational specialization: action and orientation}

The elliptic action in \cref{x_reciprocidad:iv:ellipse:section,x_reciprocidad:nuclear:7} supplies a physical specialization of the principle, with parameters constructed at their HMT stages. In its canonical coordinates,

\[
H_\eta(Q,P)=\frac{\omega}{2}
\left(e^{-\eta}Q^2+e^\eta P^2\right),\qquad
I_\eta=H_\eta/\omega.
\]

The canonical transformation \(Q=e^{\eta/2}q\), \(P=e^{-\eta/2}p\) turns \(I_\eta\) into \((q^2+p^2)/2\). Rotations of \((q,p)\), transported to \((Q,P)\), form a continuous symmetry. For the Lagrangian \(L=P\dot Q-\omega I_\eta\), its Noether charge is \(I_\eta\). On the source contour \(I_\eta=\hbar\), the oriented action is \(\oint P\,dQ=2\pi\hbar=h\), with the orientation of the Hamiltonian orbit. The proof distinguishes the symplectic and antisymplectic transformations: they can preserve the same energy while acting with different signs on oriented action. The full variational proof is presented in \cref{x_reciprocidad:nuclear:7}.

\subsection{Updating the form and the connection term}

The form matrix \(M_\eta=\operatorname{diag}(e^{\eta/2},e^{-\eta/2})\) allows development of transport between ellipses with different parameters. Fix \(J_2=\begin{pmatrix}0&-1\\1&0\end{pmatrix}\) and \(R(\theta)=\exp(-\theta J_2)\), the Hamiltonian rotation in the \((q,p)\)-plane. The transport

\[
z'=M_{\eta'}R(\delta\theta)M_\eta^{-1}z,
\qquad z=(Q,P)^\mathsf T,
\]

is symplectic and satisfies exactly \(I_{\eta'}(z')=I_\eta(z)\). The proof consists of moving both states to their normalized coordinates and using preservation of \(q^2+p^2\) under rotation. It preserves action while changing the form.

For a differentiable realization \(\eta(t)\), with angular advance \(\dot\theta=\omega\), differentiation of this transport gives

\[
\dot Q=\omega e^\eta P+\tfrac12\dot\eta Q,\qquad
\dot P=-\omega e^{-\eta}Q-\tfrac12\dot\eta P.
\]

Its Hamiltonian generator, obtained from these two equations, is

\begin{equation}
\boxed{H_{\rm trans}=\omega I_\eta+\frac{\dot\eta}{2}QP.}\label{x_reciprocidad:nuc02:eq:7}
\end{equation}

The second summand is the connection term for the varying form. With the convention \(\{f,g\}=\partial_Qf\partial_Pg-\partial_Pf\partial_Qg\),

\[
\partial_tI_\eta=\frac{\dot\eta}{2}(-e^{-\eta}Q^2+e^\eta P^2),
\quad
\left\{I_\eta,\frac{\dot\eta}{2}QP\right\}
=\frac{\dot\eta}{2}(e^{-\eta}Q^2-e^\eta P^2).
\]

The two contributions cancel, and \(dI_\eta/dt=0\). For the prescribed transport, \eqref{x_reciprocidad:nuc02:eq:7} is unique up to a function of time. The differentiable realization is a declared specialization of transport between forms; the preceding discrete result applies directly between two generated parameters. The specific TPK update retains its own rule for selecting these parameters and their times.

This gives a precise connection: an observable conserved by the complete dynamics is distributed between readout and memory through \eqref{x_reciprocidad:nuc02:eq:1}. When this observable is the generator of a continuous symmetry of the constructed action, the distributed quantity is a Noether charge. The norm, the symmetry charge, and the energy functional are related through their specific operators.

An effective bounded specialization is obtained from the 108-state calendar. Its realized Hamiltonian is \(H_{\rm clk}=\hbar\omega_0N_{108}\), with \(N_{108}=\sum_{k=0}^{107}kP_k\) and \(P_k\) its Fourier projectors. The clock action has charge \(\hbar\langle\psi,N_{108}\psi\rangle=E_{\rm clk}/\omega_0\). Transports and inclusions satisfying the hypotheses of \eqref{x_reciprocidad:nuc02:eq:1} distribute this charge between readout and complement. All operators here are finite and bounded. The elliptic action of \eqref{x_reciprocidad:nuc02:eq:7} has been proved in its classical symplectic realization; a quantization with unbounded observables additionally retains its own domain conditions.
